\chapter*{What This Booklet Is}
\addcontentsline{toc}{chapter}{What This Booklet Is}
\markboth{What This Booklet Is}{What This Booklet Is}

\textit{Mathesis} isn't a finished textbook. It's a working record of one
person's attempt to understand mathematics from its roots --- logic and
sets on up --- documented booklet by booklet as the understanding is
actually built, not summarized afterward from a position of already
knowing it.

That has one direct consequence for what you're holding: some chapters
here are fully written, and some are still outlines --- a title, a
one-paragraph statement of what the chapter needs to establish and why
it matters, and a note on what it depends on. An outline chapter says so
plainly, right where it appears. Nothing is dressed up as finished when
it isn't.

A few conventions carry over from booklet to booklet:

\begin{notebox}
Passages set off like this one flag something the main text leans on
without stopping to justify: a term worth pinning down, a distinction
easy to blur, or an aside that matters more than its size suggests.
\end{notebox}

\begin{tipbox}
This style marks a practical suggestion --- a different way to read a
passage, or a pointer toward something worth chasing further.
\end{tipbox}

Each booklet corresponds to one part of the larger project and can be
read on its own; where a chapter depends on something from an earlier
booklet, that dependency is named directly rather than assumed. Beyond
that, there's no fixed order and no hidden prerequisite this booklet
doesn't tell you about.
